\documentclass[a4paper,10pt]{report}\input{../0.Préambule/Préambule_cours_prof}\begin{document}

\pagestyle{empty}
\newpage
\noindent NOM, Prénom et classe : \dotfill

\section*{Égalité de Pythagore - Nombres rationnels \hspace{3cm} Sujet A}
\addcontentsline{toc}{section}{Égalité de Pythagore - Nombres rationnels \hspace{3cm} Sujet A}

\noindent \textit{L'usage de la calculatrice est \textbf{autorisée}.} \hfill \textit{Durée : 55 minutes}

\paragraph{Exercice 1}: Question de cours\hfill \textbf{2 points}

\medskip
Donner l'égalité de Pythagore dans un triangle $TUC$ rectangle en $U$.

\paragraph{Exercice 2}: \textit{Quotient de relatifs} \hfill \textbf{2 points}

\begin{enumerate}
\item Trouve deux nombres relatifs dont le quotient est compris entre $0$ et $1$ et la somme est négative.
%\item Trouve deux nombres relatifs dont le quotient est compris entre $0$ et $-1$ et la somme est positive.
%\item Trouve deux nombres relatifs dont le quotient et la somme sont positifs.
\item Trouve deux nombres relatifs dont le quotient et la somme dont négatifs.
\end{enumerate}

\paragraph{Exercice 3}: \textit{Arrondis de nombres}\hfill \textbf{3 points}

\begin{enumerate}
\item En utilisant la calculatrice, donne une valeur approchée arrondi au centième
\begin{enumerate}
\begin{tabularx}{\linewidth}{*{3}{>{\arraybackslash}X}}
\item $ (-1) \div 3$ &
%\item $ (-5) \div (-11)$ &
\item $ 1,3 \div 0,7$ &
%\item $\dfrac{11}{-19}$ &
\item $\dfrac{3}{5}$ \\
%\item $\dfrac{-1}{-7}$ \\
\end{tabularx}
\end{enumerate}

\item En utilisant la calculatrice, donne une valeur approchée troncaturée au millième
\begin{enumerate}
\begin{tabularx}{\linewidth}{*{3}{>{\arraybackslash}X}}
%\item $-\dfrac{-53}{16}$ &
\item $-\dfrac{-17}{-47}$ &
%\item $-\dfrac{-1,7}{-0,7}$ &
\item $ 2,9 \div (-6)$ &
%\item $ 47 \div (-23)$ &
\item $ -9,5 \div 7$ \\
\end{tabularx}
\end{enumerate}
\end{enumerate}

\paragraph{Exercice 4}: \hfill \textbf{4 points}

\noindent \textit{La rédaction sera évaluée au même titre que les résultats, on prendra donc soin de justifier chacun de nos calculs.}

\medskip
$BUS$ est un triangle rectangle en $B$ tel que $BS= 5$cm et $SU=9$cm

\medskip
\noindent Calculer la longueur $UB$, arrondie au millimètre près.

\paragraph{Exercice 4}: \hfill \textbf{5 points}

\noindent \textit{C'est un exercice de modélisation, toute trace de recherche, même non aboutie pourra être valorisée.}

\medskip
En l'an 3019 T.A., au début de la Guerre de l'Anneau, suite aux batailles des Gués de l'Isen, le roi Théoden de Rohan conduisit une armée de renfort mais fit retraite vers le Gouffre de Helm pour se réfugier de l'attaque imminente de l'armée de Saroumane, composée de redoutable orcs modifiés, appellé Uruk-hai.

\begin{quotation}
Le Gouffre était protégé par le long Mur du Gouffre, un solide mur de pierre, percé afin de laisser passer les eaux de la Rivière du Gouffre qui permettait de s'approvisionner en eau douce en situation de siège. Le mur mesure $20$ pieds de haut et est suffisamment large pour permettre à quatre hommes de se tenir côte-à-côte.
\end{quotation}

\begin{flushright}
\textit{Le Seigneur des anneaux, Livre IV, chapitre 7 « La Gorge de Helm »}
\end{flushright}

 Afin d'envahir la place fortifiée, ces Uruk-hai durent franchir, avec une échelle, ce mur devant lequel se trouve un fossé rempli d'eau, d'une largeur de $2$m. 

\medskip
\noindent Sachant que $1$ pied mesure environ $30$cm, quelle doit être la longueur minimum de l'échelle qu'ils ont du construire ?

\paragraph{Exercice 6}: \hfill \textbf{3 points}

\medskip
\noindent \textit{Cet exercice est plus compliqué que les autres et ne rapporte pas beaucoup de points, je vous conseille de le traiter en dernier une fois tous les autres terminés.} 

\noindent \textit{La rédaction sera évaluée au même titre que les résultats, on prendra donc soin de justifier chacun de nos calculs.} 

\medskip
$TSF$ est un triangle équilatéral de coté $4$cm.
\begin{enumerate}
\item Calcule la longueur d'une hauteur de ce triangle.
\item Déduis-en l'aire de ce triangle.
\end{enumerate}

\newpage
\noindent NOM, Prénom et classe : \dotfill

\section*{Égalité de Pythagore - Nombres rationnels \hspace{3cm} Sujet B}
\addcontentsline{toc}{section}{Égalité de Pythagore - Nombres rationnels \hspace{3cm} Sujet B}

\noindent \textit{L'usage de la calculatrice est \textbf{autorisée}.} \hfill \textit{Durée : 55 minutes}

\paragraph{Exercice 1}: Question de cours\hfill \textbf{2 points}

\medskip
Donner l'égalité de Pythagore dans un triangle $TUC$ rectangle en $T$.

\paragraph{Exercice 2}: \textit{Quotient de relatifs} \hfill \textbf{2 points}

\begin{enumerate}
%\item Trouve deux nombres relatifs dont le quotient est compris entre $0$ et $1$ et la somme est négative.
\item Trouve deux nombres relatifs dont le quotient est compris entre $0$ et $-1$ et la somme est positive.
\item Trouve deux nombres relatifs dont le quotient et la somme sont positifs.
%\item Trouve deux nombres relatifs dont le quotient et la somme dont négatifs.
\end{enumerate}

\paragraph{Exercice 3}: \textit{Arrondis de nombres} \hfill \textbf{3 points}

\begin{enumerate}
\item En utilisant la calculatrice, donne une valeur approchée arrondi au centième
\begin{enumerate}
\begin{tabularx}{\linewidth}{*{3}{>{\arraybackslash}X}}
%\item $ (-1) \div 3$ &
\item $ (-5) \div (-11)$ &
%\item $ 1,3 \div 0,7$ &
\item $\dfrac{11}{-19}$ &
%\item $\dfrac{3}{5}$ &
\item $\dfrac{-1}{-7}$ \\
\end{tabularx}
\end{enumerate}

\item En utilisant la calculatrice, donne une valeur approchée troncaturée au millième
\begin{enumerate}
\begin{tabularx}{\linewidth}{*{3}{>{\arraybackslash}X}}
\item $-\dfrac{-53}{16}$ &
%\item $\dfrac{-17}{-47}$ &
\item $\dfrac{-1,7}{-0,7}$ &
%\item $ 2,9 \div (-6)$ &
\item $ 47 \div (-23)$ \\
%\item $ -9,5 \div 7$ \\
\end{tabularx}
\end{enumerate}
\end{enumerate}

\paragraph{Exercice 4}: \hfill \textbf{5 points}

\noindent \textit{La rédaction sera évaluée au même titre que les résultats, on prendra donc soin de justifier chacun de nos calculs.}

\medskip
$BUS$ est un triangle rectangle en $B$ tel que $BS= 5$cm et $SU=7$cm

\medskip
\noindent Calculer la longueur $UB$, arrondie au millimètre près.

\paragraph{Exercice 4}: \hfill \textbf{5 points}

\noindent \textit{C'est un exercice de modélisation, toute trace de recherche, même non aboutie pourra être valorisée.}

\medskip
En l'an 3019 T.A., au début de la Guerre de l'Anneau, suite aux batailles des Gués de l'Isen, le roi Théoden de Rohan conduisit une armée de renfort mais fit retraite vers le Gouffre de Helm pour se réfugier de l'attaque imminente de l'armée de Saroumane, composée de redoutable orcs modifiés, appellé Uruk-hai.

\begin{quotation}
Le Gouffre était protégé par le long Mur du Gouffre, un solide mur de pierre, percé afin de laisser passer les eaux de la Rivière du Gouffre qui permettait de s'approvisionner en eau douce en situation de siège. Le mur mesure $20$ pieds de haut et est suffisamment large pour permettre à quatre hommes de se tenir côte-à-côte.
\end{quotation}

\begin{flushright}
\textit{Le Seigneur des anneaux, Livre IV, chapitre 7 « La Gorge de Helm »}
\end{flushright}

 Afin d'envahir la place fortifiée, ces Uruk-hai durent franchir, avec une échelle, ce mur devant lequel se trouve un fossé rempli d'eau, d'une largeur de $1,5$m. 

\medskip
\noindent Sachant que $1$ pied mesure environ $30$cm, quelle doit être la longueur minimum de l'échelle qu'ils ont du construire ?

\paragraph{Exercice 6}: \hfill \textbf{3 points}

\medskip
\noindent \textit{Cet exercice est plus compliqué que les autres et ne rapporte pas beaucoup de points, je vous conseille de le traiter en dernier une fois tous les autres terminés.} 

\noindent \textit{La rédaction sera évaluée au même titre que les résultats, on prendra donc soin de justifier chacun de nos calculs.} 

\medskip
$TSF$ est un triangle équilatéral de coté $6$cm.
\begin{enumerate}
\item Calcule la longueur d'une hauteur de ce triangle.
\item Déduis-en l'aire de ce triangle.
\end{enumerate}

\end{document}
